% -----------------------------------------------------------------------------
% This file is automatically generated.
% Do not edit manually.
% Source: notebooks/support/probability-distributions/support.Rmd
% -----------------------------------------------------------------------------


\noindent Hypothesis testing is done with the function \ttblue{qf}, where we use the argument \ttblue{lower.tail = FALSE} to indicate the direction of the test.

\par\addvspace{\topsep}
\begingroup
\fvset{listparameters={%
  \setlength{\topsep}{0pt}%
  \setlength{\partopsep}{0pt}%
  \setlength{\parsep}{0pt}%
  \setlength{\itemsep}{0pt}%
}}
\begin{Verbatim}[breaklines=true,formatcom=\color{ada_blue}]
(f_crit <- qf(p = 0.05, df1 = 25, df2 = 23, lower.tail = FALSE))
\end{Verbatim}
\begin{Verbatim}[breaklines=true,formatcom=\color{ada_light_blue}]
[1] 1.996271
\end{Verbatim}
\endgroup
\par\addvspace{\topsep}

\noindent In the example above, the critical value of the $F$ statistic is 1.996271. Depending on whether the value of the test statistic is greater or smaller than this value, we then reject the null hypothesis or we do not.

\noindent The probability of obtaining this particular value of the test statistic is calculated with the \ttblue{pf} function.
